% Бүлэг 1: Онолын судалгаа

\chapter{Хөгжим боловсруулалт болон хиймэл оюун ухааны онолын судалгаа}
\label{Chapter1}

%-------------------------------------------------------------------------------
\section{Дижитал хөгжим боловсруулалтын өнөөгийн байдал}

Орчин үеийн хөгжмийн үйлдвэрлэлд Дижитал Аудио Ажлын Станц буюу DAW (Digital Audio Workstation) системүүд голлох үүргийг гүйцэтгэж байна. Гэсэн хэдий ч дууны инженерчлэл, тэр дундаа миксинг (mixing) болон мастеринг (mastering) үйл ажиллагаа нь гүнзгий түвшний физик (акустик) болон математик (дохио боловсруулалт) мэдлэг шаарддаг.

Эхлэн суралцагчдад тулгардаг үндсэн хүндрэлүүдийг дараах байдлаар ангилж болно:
\begin{itemize}
    \item \textbf{Техникийн нарийн төвөгтэй байдал:} EQ (Equalization), Dynamic Range Compression зэрэг хэрэгслүүдийн параметрийг зөв тохируулах нь туршлагагүй хэрэглэгчийн хувьд субъектив алдаа гаргахад хүргэдэг.
    \item \textbf{Мэдээллийн хэт ачаалал:} Маш олон төрлийн VST (Virtual Studio Technology) плагинууд болон аудио дээжүүдээс (samples) сонголт хийх "Decision Fatigue" буюу шийдвэр гаргалтын тушилт үүсдэг.
    \item \textbf{Стандарт бус сургалтын материал:} Интернет орчинд байгаа зааварчилгаанууд нь баталгаажаагүй, эсвэл хэт ерөнхий байдаг нь суралцах явцыг удаашруулдаг.
\end{itemize}

%-------------------------------------------------------------------------------
\section{Аудио дохио боловсруулалт ба Спектраль анализ}

Дижитал аудио дохио нь хугацааны хамаарал бүхий дискрет утгуудын дараалал юм. Түүнийг хиймэл оюун ухаанаар боловсруулахын тулд хамгийн түрүүнд давтамжийн муж (Frequency Domain) руу хөрвүүлэх шаардлагатай болдог.



Богино хугацааны Фурье хувиргалт (Short-Time Fourier Transform - STFT) нь аудио дохиог Spectrogram буюу визуал дүрслэл болгон хувиргадаг. Энэхүү дүрслэл нь $x$-тэнхлэгт хугацаа, $y$-тэнхлэгт давтамж, харин өнгөний эрчмээр далайцыг (amplitude) илэрхийлдэг. Энэ нь CNN загвар аудиог "зураг" мэтээр хүлээн авч, гүнзгий анализ хийх суурь нөхцөл болдог.

%-------------------------------------------------------------------------------
\section{Convolutional Neural Networks (CNN) ба Аудио танилт}

CNN (Хувиргалтын мэдрэлийн сүлжээ) нь өгөгдлийн орон зайн хэв шинжийг (spatial features) илрүүлэхдээ гарамгай архитектур юм. Дууны миксийн чанарыг үнэлэхдээ CNN-ийг дараах байдлаар ашиглана:
\begin{enumerate}
    \item \textbf{Feature Extraction:} Олон давхаргат хувиргалтын шүүлтүүрүүд ашиглан аудио спектрограмм дахь давтамжийн тэнцвэрт байдал, шуугианы түвшинг илрүүлнэ.
    \item \textbf{Classification/Regression:} Сургагдсан загвар нь тухайн аудиог мэргэжлийн түвшний ("сайн") эсвэл алдаатай ("муу") микс мөн эсэхийг магадлалын онолоор тодорхойлно.
\end{enumerate}



%-------------------------------------------------------------------------------
\section{Генератив загварууд: MusicVAE}

Хөгжим зохиох үйл явцад хүн-компьютерийн харилцан ажиллагааг (HCI) нэмэгдүүлэхийн тулд Variational Autoencoder (VAE) загварыг ашигладаг. MusicVAE нь хөгжмийн дарааллыг латент орон зайд (Latent Space) кодлодог бөгөөд энэ нь:
\begin{itemize}
    \item Хоёр өөр аяны хооронд шилжилт (Interpolation) хийх;
    \item Одоо байгаа аян дээр тулгуурлан шинэ хувилбарууд (Reconstruction) үүсгэх боломжийг олгодог.
\end{itemize}
Энэ нь хэрэглэгчид бүтээлч байдлын гацалтаас гарахад туслах AI-ментор болох юм.

%-------------------------------------------------------------------------------
\section{LLM ба RAG технологийн хэрэглээ}

Орчин үеийн Том хэлний загварууд (LLM) хөгжмийн техникийн зөвлөгөө өгөх чадвартай боловч худал мэдээлэл (Hallucination) үүсгэх эрсдэлтэй. Үүнийг шийдвэрлэхийн тулд Мэдээлэл хайлтаар баяжуулсан үүсгэлт буюу RAG (Retrieval-Augmented Generation) аргыг ашиглана.



RAG нь хэрэглэгчийн асуултад хариулахын өмнө мэргэжлийн продюсеруудын баримт бичиг, FL Studio-ийн албан ёсны гарын авлага бүхий Vector Database-ээс хайлт хийж, зөвхөн баталгаажсан мэдээлэлд суурилсан хариултыг боловсруулдаг.

%-------------------------------------------------------------------------------
\section{Бүлгийн дүгнэлт}

Нэгдүгээр бүлгээр дижитал аудио боловсруулалтын математик үндэс болон хиймэл оюун ухааны орчин үеийн загваруудыг (CNN, MusicVAE, RAG) судлан үзлээ. Эдгээр онолын судалгаа нь дууны миксийн чанарыг AI-аар үнэлэх, хэрэглэгчид ухаалаг зөвлөгөө өгөх нэгдсэн системийг хөгжүүлэх бодит боломжтойг баталж байна. Дараагийн бүлэгт эдгээр технологид суурилсан системийн архитектур, загварчлалыг авч үзэх болно.